% Basisapparatuur: a \subchapter-level (H2) heading grouped under the single
% "Inleiding" chapter-level heading in main.tex.
\subchapter{Basisapparatuur}

\lettrine{G}{een} verplichte checklist, maar wat ik zelf in de keuken heb
staan. Met deze apparatuur erbij kun je bij vrijwel elk recept in dit boek
gewoon beginnen, van een simpele hapjespan tot de dingen die maar af en toe
uit de kast komen.

\blockrule{Pannen \& potten}
\begin{ingredients}
\ingb{Braadpan, ovenvast}
\ingbnote{Hapjespan}{ (28 cm), met deksel}
\ingb{Koekenpan (24 cm)}
\ingbnote{Koekenpan, klein}{ (20 cm)}
\ingb{Soeppan (5 liter)}
\ingb{Steelpannetje (16 cm)}
\ingb{Wokpan (30 cm)}
\end{ingredients}

\blockrule{Bakvormen \& ovenschalen}
\begin{ingredients}
\ingb{Bakplaat (ca.\ 30 × 40 cm)}
\ingbnote{Braadslee}{ (ca.\ 30 × 40 cm), voor ovenfriet en grotere braadstukken}
\ingb{Cakevorm (25 cm)}
\ingb{Diepe ovenschaal (ca.\ 32 × 20 × 8 cm)}
\ingb{Muffinvorm (12 holtes)}
\ingb{Pizzasteen, rond (ca.\ 30 cm)}
\ingbnote{Springvorm}{ (24 cm), doet ook dienst als quichevorm}
\ingb{Taartvorm, rond (18 cm)}
\ingb{Vierkante bakvorm (20 × 20 cm)}
\end{ingredients}

\clearpage
\blockrule{Kleine apparatuur}
\tip{Geen keukenmachine of blender in huis: de staafmixer dekt fijnmalen en
pureren, direct in de pan of in een losse kom, en de standmixer neemt het
kneden over. Bij gekookte aardappelpuree werkt een staafmixer juist
averechts, die maakt de puree plakkerig in plaats van luchtig, gebruik dan
een pureeknijper.}
\begin{ingredients}
\ingbnote{Pastamachine}{, optioneel, een deegroller kan ook}
\ingbnote{Staafmixer}{, voor het pureren en fijnmalen, in de pan of in een losse kom}
\ingbnote{Standmixer}{ (b.v.\ KitchenAid Artisan), voor het kneden van brood- en pizzadeeg}
\end{ingredients}

\blockrule{Handgereedschap}
\begin{ingredients}
\ingb{Deegroller}
\ingbnote{Fijne zeef}{, voor een extra gladde soep}
\ingb{Garde}
\ingbnote{Garneerspuit}{ (Leifheit, eenhandig), met 6 spuitmondjes}
\ingbnote{Kaasschaaf}{, ook voor flinterdunne groente}
\ingbnote{Pureeknijper}{ (Joseph Joseph Helix), voor aardappelpuree}
\ingbnote{Rasp}{, voor citroenrasp en kaas}
\ingb{Schuimspaan}
\ingb{Vergiet}
\end{ingredients}

\blockrule{Meetinstrumenten}
\begin{ingredients}
\ingbnote{Digitale kernthermometer}{, voor vlees en gefrituurd eten}
\ingb{Maatbeker (1 liter)}
\ingbnote{Weegschaal}{ (KitchenAid, dubbel platform), met een vlak tot 5 kg voor grotere hoeveelheden en een precisievlak tot 0,01 g voor gist, kruiden en zout}
\end{ingredients}

\vspace{0.6em}
{\footnotesize\itshape\color{inkfaint}Dit is wat er bij ons in de keuken
staat, geen verplichte lijst. Niet alles heb je voor elk recept nodig, en
het meeste is prima te vervangen door iets wat je al hebt.\par}
